\chapter{NORMED LINEAR SPACES}

\section{Norms}
In the world of analysis the predominant denizens are function spaces, vector spaces of
real or complex valued functions. To be of interest to an analyst such a space should
come equipped with a topology.  Often the topology is a metric topology, which in turn
frequently comes from a norm (see proposition~\ref{0001503}).

\begin{exam}\label{C023125}  For $x = (x_1,\dots, x_n) \in \K^n$ let $\norm{x} =
\left(\sum_{k=1}^n {\abs{x_k}}^2\right)^{1/2}$.  This is the
 \index{usual!norm!on $\K^n$}%
 \index{norm!usual!on $\K^n$}%
\df{usual norm} (or
 \index{Euclidean!norm!on $\K^n$}%
 \index{norm!Euclidean!on $\K^n$}%
 \index{K@$\K^n$!as a normed linear space}%
 \index{normed!linear space!$\K^n$ as a}%
\df{Euclidean norm}) on $\K^n$; unless the contrary is explicitly stated, $\K^n$ when
regarded as a normed linear space will always be assumed to possess this norm.  It is
clear that this norm induces the usual (Euclidean) metric on~$\K^n$.
\end{exam}

\begin{exam}\label{C023127}  For $x = (x_1,\dots,x_n) \in \K^n$ let
 \index{<unopn@$\norm{\hphantom{x}}_1$ ($1$-norm)}%
$\norm{x}_1 = \sum_{k=1}^n \abs{x_k}$.  The function $x \mapsto \norm{x}_1$ is easily
seen to be a norm on~$\K^n$. It is sometimes called
 \index{norm!$1$- (on $\K^n$)}%
 \index{one@$1$-norm!on $\K^n$}%
 \index{K@$\K^n$!as a normed linear space}%
 \index{normed!linear space!$\K^n$ as a}%
the \df{$1$-norm} on~$\K^n$.  It induces the so-called \emph{taxicab metric} on~$\R^2$.
\end{exam}

Of the next six examples the last is the most general.  Notice that all the others are special cases of it
or subspaces of special cases.

\begin{exam}\label{C023129}  For $x = (x_1,\dots,x_n) \in \K^n$ let
 \index{<unopn@$\norm{\hphantom{x}}_u$ (uniform norm)}%
 \index{<unopn@$\norm{\hphantom{x}}_\infty$ (uniform norm)}%
$\norm{x}_\infty = \max\{\,\abs{x_k} \colon 1 \le k \le n\}$.  This defines a norm on $\K^n$;
it is the
 \index{uniform!norm!on~$\K^n$}%
 \index{norm!uniform!on $\K^n$}%
 \index{K@$\K^n$!as a normed linear space}%
 \index{normed!linear space!$\K^n$ as a}%
\df{uniform norm} on~$\K^n$.  It induces the uniform metric on~$\K^n$.  An alternative notation for
$\norm{x}_\infty$ is $\norm{x}_U$.
\end{exam}

\begin{exam}\label{C0231301} The set $l_\infty$ of all bounded sequences of complex numbers is
clearly a vector space under pointwise operations of addition and scalar multiplication.
For $x = (x_1,x_2,\dots,) \in l_\infty$ let
 \index{<unopn@$\norm{\hphantom{x}}_\infty$ (uniform norm)}%
$\norm{x}_\infty = \sup\{\,\abs{x_k} \colon 1 \le k\}$.  This defines a norm on
$l_\infty$; it is the
 \index{uniform!norm!on~$l_\infty$}%
 \index{norm!uniform!on $l_\infty$}%
 \index{l@$l_\infty$!as a normed linear space}%
 \index{l@$l_\infty$!bounded sequences}%
 \index{normed!linear space!$l_\infty$ as a}%
\df{uniform norm} on~$l_\infty$.
\end{exam}

\begin{exam}\label{C0231302} A subspace of~\ref{C0231301} is the space $c$ of all convergent 
 \index{c@$c$!as a normed linear space}%
 \index{c@$c$!convergent sequences}%
 \index{normed!linear space!$c$ as a}%
sequences of complex numbers.  Under the uniform norm it is a normed linear space.
\end{exam}

\begin{exam}\label{C0231303} Contained in $c$ is another interesting normed linear space $\sbsb c0$
 \index{c@$c_0$!sequences which converge to zero}%
 \index{c@$c_0$!as a normed linear space}%
 \index{normed!linear space!$\sbsb c)$ as a}%
the space of all sequences of complex numbers which converge to zero (again with the uniform norm).
\end{exam}

\begin{exam}\label{C023131}  Let $S$ be a nonempty set.  If $f$ is a bounded $K$-valued function on $S$, let
  \[ \norm f_u := \sup\{\abs{f(x)} \colon x \in S\}\,. \]
This is a norm on the vector space $\fml B(S)$ of all bounded complex valued functions on $S$ and is called the
 \index{<unopn@$\norm{\hphantom{x}}_u$ (uniform norm)}%
 \index{uniform!norm!on~$\fml B(S)$}%
 \index{norm!uniform!on~$\fml B(S)$}%
 \index{B@$\fml B(S)$!bounded functions on~$S$}%
 \index{B@$\fml B(S)$!as a normed linear space}%
 \index{normed!linear space!$\fml B(S)$ as a}%
\df{uniform norm}.  Clearly this norm gives rise to the uniform metric on~$\fml B(S)$.
Notice that examples~\ref{C023129} and~\ref{C0231301} are
special cases of this one. (Let $S = \N_n := \{1,2,\dots,n\}$ or $S = \N$.)
\end{exam}











\begin{exam}\label{C023133} It is easy to make a substantial generalization of the preceding
example by replacing the field $\K$ by an arbitrary normed linear space. Suppose then
that $S$ is a nonempty set and $V$ is a normed linear space. For $f$ in the vector space $\fml B(S,V)$ of all bounded $V$-valued functions on~$S$, let
  \[ \norm f_u := \sup\{\norm{f(x)} \colon x \in S\}\,. \]
Then this is a norm and it is called the
 \index{<unopn@$\norm{\hphantom{x}}_u$ (uniform norm)}%
 \index{uniform!norm!on~$\fml B(S,V)$}%
 \index{norm!uniform!on~$\fml B(S,V)$}%
 \index{B@$\fml B(S,V)$!Bounded $V$-valued functions on~$S$}%
 \index{B@$\fml B(S,V)$!as a normed linear space}%
 \index{normed!linear space!$\fml B(S,V)$ as a}%
\df{uniform norm} on $\fml B(S,V)$.   (Why is the word ``easy'' in the first sentence of this example appropriate?)
\end{exam}

\begin{defn} Let $S$ be a set, $V$ be a normed linear space, and
 \index{F@$\fml F(S,V)$ ($V$-valued functions on~$S$)}%
$\fml F(S,V)$ be the vector space of all $V$-valued functions on~$S$.
Consider a sequence $(f_n)$ of functions in $\fml F(S,V)$.  If there is a function $g$ in $\fml F(S,V)$ such that
   \[ \sup\{\norm{f_n(x) - g(x)}\colon x \in S\} \sto 0 \quad \text{as $n \sto \infty$,} \]
then we say that the sequence $(f_n)$
 \index{<arrowto@$f_n \sto g$ (unif) (uniform convergence)}%
 \index{uniform!convergence}%
 \index{convergence!uniform}%
\df{converges uniformly} to~$g$ and write $f_n \sto g \text{\,(unif)}$.  The function $g$
is the
 \index{uniform!limit}%
 \index{limit!uniform}%
\df{uniform limit} of the sequence $(f_n)$.  Notice that if $g$ and all the $f_n$'s
belong to $\fml B(S,V)$, then uniform convergence of $(f_n)$ to $g$ is just convergence
of $(f_n)$ to $g$ with respect to the uniform metric.
\end{defn}

There are many ways in which sequences of functions converge.  Arguably the two most
common modes of convergence are uniform convergence, which we have just discussed, and
pointwise convergence.

\begin{defn}  Let $S$ be a set, $V$ be a normed linear space, and $(f_n)$ be a sequence in
$\fml F(S,V)$.  If there is a function $g$ such that
  \[ f_n(x) \sto g(x) \qquad  \text{for all $x \in S$}, \]
then $(fn)$
 \index{pointwise!convergence}%
 \index{convergence!pointwise}%
\df{converges pointwise} to~$g$.  In this case we write
  \[ f_n \sto g \text{\,(ptws)}. \]
The function $g$ is the
 \index{pointwise!limit}%
 \index{limit!pointwise}%
 \index{<arrowto@$f_n \sto g$ (ptws) (pointwise convergence)}%
\df{pointwise limit} of the~$f_n$'s.

If $(f_n)$ is an increasing sequence of real (or extended real) valued functions and $f_n
\sto g \text{\,(ptws)}$, we write
 \index{<arrowmonoup@$f_n \uparrow g$ (ptws) (pointwise monotone convergence)}%
$f_n \uparrow g \text{\,(ptws)}$. And if $(f_n)$ is decreasing and has $g$ as a pointwise
limit, we write
 \index{<arrowmonodown@$f_n \downarrow g$ (ptws) (pointwise monotone convergence)}%
$f_n \downarrow g \text{\,(ptws)}$.
\end{defn}

The most obvious connection, familiar from any real analysis course, between these two types of convergence is that uniform
convergence implies pointwise convergence.

\begin{prop}\label{C023141} Let $S$ be a set and $V$ be a normed linear space.  If a sequence $(f_n)$ in $\fml F(S,V)$
converges uniformly to a function $g$ in $\fml F(S,V)$, then $(f_n)$ converges pointwise to~$g$.
\end{prop}

The converse is not true.

\begin{exam}  For each $n \in \N$ let $f_n\colon [0,1] \sto \R\colon x \mapsto x^n$. Then the
sequence $(f_n)$ converges pointwise on $[0,1]$, but not uniformly.
\end{exam}

\begin{exam} For each $n \in \N$ and each $x \in \R^+$ let $f_n(x) = \frac1n\, x$.  Then on
each interval of the form $[0,a]$ where $a > 0$ the sequence $(f_n)$ converges uniformly
to the constant function~$0$.  On the interval $[0,\infty)$ it converges pointwise to
$0$, but not uniformly.
\end{exam}

Recall also from real analysis that a uniform limit of bounded functions must itself be bounded.

\begin{prop}\label{C023151} Let $S$ be a set, $V$ be a normed linear space, and $(f_n)$ be a
sequence in $\fml B(S,V)$ and $g$ be a member of $\fml F(S,V)$.  If $f_n \sto g\text{\,(unif)}$,
then $g$ is bounded.
\end{prop}

And a uniform limit of continuous functions is continuous.

\begin{prop}\label{C023154} Let $X$ be a topological space, $V$ be a normed linear space, and
$(f_n)$ be a sequence of continuous $V$-valued functions on~$X$ and $g$ be a member of $\fml F(X,V)$.  If $f_n
\sto g\text{\,(unif)}$, then $g$ is continuous.
\end{prop}

\begin{prop}\label{C023163} If $x$ and $y$ are elements of a vector space $V$ equipped with a
norm $\norm{\hphantom{x}}$, then
  \[ \bigabs{\strut\,\norm{x} - \norm{y}\,} \le \norm{x - y}\,. \]
\end{prop}

\begin{cor} The norm on a normed linear space is a uniformly continuous function.
 \index{continuity!of norms}%
 \index{norm!continuity of}%
\end{cor}

\begin{notn} Let $M$ be a metric space, $a \in M$, and $r > 0$.  We denote $\{x \in M\colon d(x,a) < r\}$, the
 \index{B@$B_r(a)$, $B_r$ (open ball of radius $r$)}%
 \index{open!ball}%
 \index{ball!open}%
open ball of radius $r$ about $a$, by $B_r(a)$.  Similarly, we denote $\{x \in M\colon d(x,a) \le r\}$, the
 \index{C@$C_r(a)$, $C_r$ (closed ball of radius $r$)}%
 \index{closed!ball}%
 \index{ball!closed}%
\df{closed ball} of radius $r$ about $a$, by $C_r(a)$ and $\{x \in M\colon d(x,a) = r\}$, the
 \index{S@$S_r(a)$, $S_r$ (sphere of radius $r$)}%
 \index{sphere}%
\df{sphere} of radius $r$ about $a$, by $S_r(a)$.  We will use $B_r$, $C_r$, and $S_r$ as abbreviations for $B_r(0)$, $C_r(0)$,
and $S_r(0)$, respectively.
\end{notn}

\begin{prop} If $V$ is a normed linear space, if $x \in V$, and if $r$, $s > 0$, then
 \begin{enumerate}[label=\rm{(\alph*)}]
   \item $B_r(\vc 0) = -B_r(\vc 0)$;
   \item $B_{rs}(\vc 0) = rB_s(\vc 0)$;
   \item $x + B_r(\vc 0) = B_r(x)$; and
   \item $B_r(\vc 0) + B_r(\vc 0) = 2\,B_r(\vc 0)$.  (Is it true in general that
$A + A = 2A$ when $A$ is a subset of a vector space?)
 \end{enumerate}
\end{prop}

\begin{defn}\label{C017314} If $a$ and $b$ are vectors in the vector space $V$, then the
 \index{closed!segment}%
 \index{segment, closed}%
\df{closed segment} between $a$ and $b$, denoted
 \index{<bracsqr@$[a,b]$ (closed segment in a vector space)}%
by~$[a,b]$, is $\{(1 - t)a + tb \colon 0 \le t \le 1\}$.
\end{defn}

\begin{cau}\label{C017317} Notice that there is a slight conflict between this notation for closed \emph{segments}, when
applied to the vector space $\R$ of real numbers, and the usual notation for closed \emph{intervals} in~$\R$.
In $\R$ the closed segment $[a,b]$ is the same as the closed
interval $[a,b]$ provided that $a \le b$. If $a > b$, however, the closed segment $[a,b]$
is the same as the segment $[b,a]$, it contains all numbers $c$ such that $b \le c \le
a$, whereas the closed interval $[a,b]$ is empty.
\end{cau}

\begin{defn} A subset $C$ of a vector space $V$ is
 \index{convex!set}%
\df{convex} if the closed segment $[a,b]$ is contained in $C$ whenever $a$, $b \in C$.
\end{defn}

\begin{exam}\label{C023169}  In a normed linear space every open ball is a convex set.  And so is every closed ball.
\end{exam}

\begin{prop}\label{prop_intrsctn_convex} The intersection of a family of convex subsets of a vector space is convex.
\end{prop}

\begin{defn} Let $V$ be a vector space. Recall that a \emph{linear combination} of a finite set $\{x_1, \dots, x_n\}$
of vectors in $V$ is a vector of the form $\sum_{k=1}^n \alpha_k x_k$ where $\alpha_1, \dots, \alpha_n \in \R$.
If $\alpha_1 = \alpha_2 = \dots = \alpha_n = 0$, then the linear combination is \emph{trivial}; if at least one
$\alpha_k$ is different from zero, the linear combination is \emph{nontrivial}.  A linear combination
$\sum_{k=1}^n\alpha_k x_k$ of the vectors $x_1, \dots, x_n$ is a
 \index{convex!combination}%
 \index{combination!convex}%
\df{convex combination} if $\alpha_k \ge 0$ for each $k$ ($1 \le k \le n$) and if $\sum_{k=1}^n \alpha_k = 1$.
\end{defn}

\begin{defn}\label{smplx005def} Let $A$ be a nonempty subset of a vector space~$V$.  We define the
 \index{convex!hull}%
 \index{hull!convex}%
 \index{co@$\co A$ (convex hull of~$A$)}%
\df{convex hull} of $A$, denoted by $\co A$, to be the smallest convex subset of $V$ which contain~$A$.
\end{defn}

The preceding definition should provoke instant suspicion.  How do we know there \emph{is} a `smallest' convex subset of~$V$
which contains~$A$?  One fairly obvious candidate (let's call it the `abstract' one) for such a set is the intersection
of \emph{all} the convex subsets of $V$ which contain~$A$.  (Of course, this might also be nonsense---\emph{unless} we know that there
is at least one convex set in $V$ which contains~$A$ and that the intersection of the family of all convex sets containing~$A$ is
itself a convex set containing~$A$.)  Another plausible candidate (we'll call it the `constructive' one) is the set of all convex
combinations of elements of~$A$.  (But \emph{is} this a convex set?  And is it the smallest one containing~$A$?)  Finally, even if these two
approaches both make sense, do they constitute equivalent definitions of \emph{convex hull}?  That is, do they both define the same set, and
does that set satisfy definition~\ref{smplx005def}?

\begin{prop}\label{smplx006} The definitions suggested in the preceding paragraph, labeled `abstract' and `constructive', make sense and
are equivalent.  Moreover, the set they describe satisfies the definition given in~\ref{smplx005def} for \emph{convex hull}.
\end{prop}

\begin{prop} If $T\colon V \sto W$ is a linear map between vector spaces and $C$ is a convex
subset of $V$, then $T^{\sto}(C)$ is a convex subset of~$W$.
\end{prop}

\begin{prop} In a normed linear space the closure of every convex set is convex.
\end{prop}

\begin{prop} Let $V$ be a normed linear space. For each $a \in V$ the map $T_a\colon V \sto V\colon x \mapsto x + a$ (called
 \index{translation}%
\df{translation} by~$a$) is a homeomorphism.
\end{prop}

\begin{cor} If $U$ is a nonempty open subset of a normed linear space $V$, then $U - U$
contains a neighborhood of~$0$.
\end{cor}

\begin{prop} If $(x_n)$ is a sequence in a normed linear space and $x_n \sto a$, then
$\D\frac1n\sum_{k=1}^nx_k~\sto~a$.
\end{prop}

In proposition~\ref{00015025} we showed how an inner product on a vector space induces a norm on that space.
It is reasonable to ask if all norms can be generated in this fashion from an inner product.  The answer
is \emph{no}.  The next proposition gives a very simple necessary and sufficient condition for a norm to
arise from an inner product.

\begin{prop} Let $V$ be a normed linear space.  There exists an inner product on $V$ which induces the
norm on $V$ if and only if the norm satisfies the \emph{parallelogram law}~\ref{parallelogram_law}.
\end{prop}

\begin{proof}[\emph{Hint for proof}]  To prove that if a norm satisfies the \emph{parallelogram law}
then it is induced by an inner product, use the equation given in the \emph{polarization identity}
(proposition~\ref{0001507}) as a definition.  Prove first that $\langle y,x \rangle = \conj{\langle x,y \rangle}$
for all $x$, $y\in V$ and that $\langle x,x \rangle > 0$ whenever $x \ne 0$.  Next, for arbitrary
$z \in V$, define a function $f\colon V \sto \C\colon x \mapsto \langle x,z \rangle$.  Prove that
   \[ f(x + y) + f(x - y) = 2f(x) \tag{$\ast$} \]
for all $x$, $y \in V$.  Use this to prove that $f(\alpha x) = \alpha f(x)$ for all $x \in V$ and $\alpha \in \R$.
(Start with $\alpha$ being a natural number.)  Then show that $f$ is additive. (If $u$ and $v$ are arbitrary
elements of $V$ let $x = u + v$ and $y = u - v$.  Use~$(\ast)$.)  Finally prove that $f(\alpha x) = \alpha f(x)$
for complex $\alpha$ by showing that $f(ix) = i\,f(x)$ for all $x \in V$.  \ns
\end{proof}

\begin{defn} Let $(X,\sfml T)$ be a topological space. A family $\sfml B \subseteq \sfml T$
is a
 \index{base!for a topology}%
\df{base} for $\sfml T$ if each member of $\sfml T$ is a union of members of~$\sfml B$.
In other words, a family $\sfml B$ of open sets is a base for $\sfml T$ if for each open
set $U$ there exists a subfamily $\sfml B'$ of $\sfml B$ such that $U = \bigcup \sfml B'$.
\end{defn}

\begin{exam}\label{xmpl_top_metric} In the plane $\R^2$ with its usual (Euclidean) topology, the family of
all open balls in the space is a base for its topology, since any open set in $\R^2$ can be expressed as a
union of open balls.  Indeed in any metric space whatever the family of open balls is a base for a topology
on the space.  This is the
 \index{topology!induced by a metric}%
 \index{induced!topology}%
\df{topology induced by the metric}.
\end{exam}

In practice (in metric spaces for example) it is often more convenient to specify a base for a topology than to specify
the topology itself.  It is important to realize, however, that there may be many
different bases for the same topology.  Once a particular base has been chosen we refer
to its members as
 \index{basic!open sets}%
 \index{open!set!basic}%
\emph{basic open sets}.

The word ``induces'' is regarded as a transitive relation.  For example, given an inner product space $V$, no one would
refer to the \emph{topology on $V$ induced by the metric induced by the norm induced by the inner product}.  One just speaks of
the \emph{topology on $V$ induced by the inner product}.  Sometimes one says ``generated by'' instead of ``induced by''.

\begin{prop} Let $V$ be a vector space with two norms $\norm{\,\cdot\,}_1$ and
$\norm{\,\cdot\,}_2$.  Let $\sfml T_k$ be the topology induced on $V$ by
$\norm{\,\cdot\,}_k$ (for $k = 1,2$). If there exists a constant $\alpha > 0$ such that
$\norm{x}_1 \le \alpha \norm{x}_2$ for every $x \in V$, then $\sfml T_1 \subseteq \sfml T_2$.
\end{prop}

\begin{defn} Two norms on a vector space $V$ are
 \index{equivalent!norms}%
 \index{norm!equivalent}%
\df{equivalent} if there exist constants $\alpha$, $\beta > 0$ such that for all $x \in
V$
   \[ \alpha\norm{x}_1 \le \norm{x}_2 \le \beta \norm{x}_1\,. \]
\end{defn}

\begin{prop}\label{C023184} If $\norm{\hphantom{x}}_1$ and $\norm{\hphantom{x}}_2$
 \index{equivalent!norms!induce identical topologies}%
are equivalent norms on a vector space~$V$, then they induce the same topology on~$V$.
\end{prop}

Proposition~\ref{C023184} gives us an easily verifiable sufficient condition for two norms
to induce identical topologies on a vector space.  Thus, if we are trying to show, for
example, that some subset of a normed linear space is open it may very well be the case
that the proof can be simplified by replacing the given norm with an equivalent one.

Similarly, suppose that we are attempting to verify that a function $f$ between two
normed linear spaces is continuous.  Since continuity is defined in terms of open sets
and equivalent norms produce exactly the same open sets (see proposition~\ref{C023184}),
we are free to replace the norms on the domain of $f$ and the codomain of $f$ with any
equivalent norms we please.  This process can sometimes simplify arguments significantly.
(This possibility of simplification, incidentally, is one major advantage of giving a
topological definition of continuity in the first place.)

\begin{defn}\label{C02319} Let $x = (x_n)$ a sequence of vectors in a normed linear space $V$.
The infinite series $\sum_{k=1}^\infty x_k$ is
 \index{convergent!series}%
 \index{series!convergent}%
\df{convergent} if there exists a vector $b \in V$ such that $\norm{b - s_n} \sto 0$ as $n
\sto \infty$, where $s_n = \sum_{k=1}^n x_k$ is the $n^{\text{th}}$ partial sum of the
sequence~$x$.  The series is
 \index{absolutely!convergent!series}%
 \index{convergent!absolutely}%
\df{absolutely convergent} if $\sum_{k=1}^\infty \norm{x_k} < \infty$.
\end{defn}

\begin{exer}\label{C023191} Let
 \index{l@$l_1$!as a normed linear space}%
 \index{l@$l_1$!absolutely summable sequences}%
 \index{normed!linear space!$l_1$ as a}%
$l_1$ be the set of all sequences $x = (x_n)$ of complex numbers such that the series $\sum x_k$ is absolutely
convergent. Make $l_1$ into a vector space with the usual pointwise definition of addition and scalar multiplication.
For every $x \in l_1$ define
 \index{<unopn@$\norm{\hphantom{x}}_1$ ($1$-norm)}%
 \index{<unopn@$\norm{\hphantom{x}}_\infty$ (uniform norm)}%
   \begin{align*}
      \norm{x}_1 &:= \sum_{k=1}^\infty \abs{x_k} \\
   \intertext{and}
      \norm{x}_\infty &:= \sup\{\abs{x_k}\colon k \in \N\}.
   \end{align*}
(These are, respectively, the
 \index{one@$1$-norm!of a sequence}%
 \index{norm!$1$- (of a sequence)}%
$1$-norm and the
 \index{uniform!norm!of a sequence}%
\emph{uniform} norm.) Show that both $\norm{\hphantom{x}}_1$ and $\norm{\hphantom{x}}_\infty$
are norms on~$l_1$. Then prove or disprove:
 \begin{enumerate}
  \item[(a)] If a sequence $(x^i)$ of vectors in the normed linear space
$\bigl(l_1, \norm{\hphantom{x}}_1\bigr)$ converges, then the sequence also converges in
$\bigl(l_1, \norm{\hphantom{x}}_\infty\bigr)$.
  \item[(b)] If a sequence $(x^i)$ of vectors in the normed linear space
$\bigl(l_1, \norm{\hphantom{x}}_\infty\bigr)$ converges, then the sequence also converges in
$\bigl(l_1, \norm{\hphantom{x}}_1\bigr)$.
 \end{enumerate}
\end{exer}









\section{Bounded Linear Maps}\label{sec_bdd_lin_maps}
Analysts work with objects having both algebraic and topological structure.  In the
preceding section we examined vector spaces endowed with a topology derived from a norm.
If these are the objects under consideration, what morphisms are likely to be of greatest
interest?  The plausible, and correct, answer is \emph{maps which preserve both
topological and algebraic structures}; that is, continuous linear maps.  The resulting
category is denoted by
 \index{category!$\cat{NLS_\infty}$ as a}%
 \index{NLSinf@$\cat{NLS_{\boldsymbol\infty}}$!the category}
$\cat{NLS_{\boldsymbol\infty}}$.  The isomorphisms in this category are, topologically,
homeomorphisms.

Here is a very useful condition which, for linear maps, turns out to be equivalent to continuity.

\begin{defn} A linear transformation $T\colon V \sto W$ between normed linear spaces is
 \index{bounded!linear map}%
 \index{linear!map!bounded}%
\df{bounded} if $T(B)$ is a bounded subset of $W$ whenever $B$ is a bounded subset
of~$V$.  In other words, a bounded linear map takes bounded sets to bounded sets.  We
denote
 \index{B@$\ofml B(V,W)$!bounded linear maps}%
by~$\ofml B(V,W)$ the family of all bounded linear transformations from $V$ into~$W$.  A
bounded linear map from a space $V$ into itself is often called a (bounded linear)
 \index{operator}%
 \index{conventions!all operators are bounded and linear}%
\df{operator}.  The family of all operators on a normed linear space $V$ is denoted
 \index{B@$\ofml B(V)$!operators}%
by~$\ofml B(V)$.  The class of normed linear spaces together with the bounded linear maps
between them constitute a category.  We verify below that this is just the
category~$\cat{NLS_{\boldsymbol\infty}}$.
\end{defn}

\begin{cau} It is extremely important to realize that a bounded linear map will \emph{not},
in general, be a bounded function in the usual sense of a \emph{bounded function} on some set
(see examples~\ref{C023131} and~\ref{C023133}).  The use of ``bounded'' in these two conflicting
senses may be unfortunate, but it is well established.
\end{cau}

\begin{exam} The linear map $T\colon \R \sto \R\colon x \mapsto 3x$ is a bounded linear map
(since it maps bounded subsets of $\R$ to bounded subsets of~$\R$), but, regarded just as
a function, $T$ is not bounded (since its range is not a bounded subset of~$\R$).
\end{exam}

The following observation my help reduce confusion.

\begin{prop} A linear transformation, unless it is the constant map that takes every vector to
zero, \emph{cannot} be a bounded function.
\end{prop}

\begin{defn}  Let $(M,d)$ and $(N,\rho)$ be metric spaces.  A function $f \colon M \sto N$ is
 \index{uniform!continuity}%
 \index{continuity!uniform}%
\df{uniformly continuous} if for every $\epsilon > 0$ there exists $\delta > 0$ such that
$\rho(f(x),f(y)) < \epsilon$ whenever  $x$, $y \in M$ and $d(x,y) < \delta$.
\end{defn}

One of the most striking aspects of linearity is that for linear maps the concepts of
continuity, continuity at a single point, and uniform continuity coalesce.  And in fact
they are exactly the same thing as boundedness.

\begin{thm}\label{C023221} Let $T \colon V \sto W$ be a linear transformation between normed
linear spaces. Then the following are equivalent:
  \begin{enumerate}[label=\rm{(\alph*)}]
    \item $T$ is bounded.
    \item The image of the closed unit ball under $T$ is bounded.
    \item $T$ is uniformly continuous on $V$.
    \item $T$ is continuous on $V$.
    \item $T$ is continuous at~$\vc 0$.
    \item There exists a number $M > 0$ such that $\norm{Tx} \le M\norm{x}$ for all $x \in V$.
  \end{enumerate}
\end{thm}

\begin{prop} Let $T\colon V \sto W$ be a bounded linear transformation between normed linear
spaces.  Then the following four numbers (exist and) are equal.
 \begin{enumerate}[label=\rm{(\alph*)}]
  \item $\sup\{\norm{Tx}\colon \norm{x} \le 1\}$
  \item $\sup\{\norm{Tx}\colon \norm{x} = 1\}$
  \item $\sup\{\norm{Tx}\,\norm{x}^{-1}\colon x \ne \vc 0\}$
  \item $\inf\{M > 0\colon \norm{Tx} \le M\norm{x}\ \text{ for all $x \in V$}\}$
 \end{enumerate}
\end{prop}

\begin{defn}\label{defn_op_norm} If $T$ is a bounded linear map, then $\norm{T}$, called the
 \index{norm!of a bounded linear map}%
 \index{<unopn@$\norm T$ (norm of a bounded linear map~$T$)}%
\df{norm} (or the
 \index{norm!operator}%
 \index{operator!norm}%
\df{operator norm}  of $T$, is defined to be any one of the four expressions in the previous exercise.
\end{defn}

\begin{exam}\label{exam_op_norm_norm} If $V$ and $W$ are normed linear spaces, then the function
   \[ \norm{\hphantom{T}}\colon \ofml B(V,W) \sto \R   \colon T \mapsto \norm{T} \]
is, in fact, a norm.
\end{exam}

\begin{exam}\label{C023228} The family $\ofml B(V,W)$ of all bounded linear maps
 \index{B@$\ofml B(V,W)$!as a normed linear space}%
 \index{normed!linear space!$\ofml B(V,W)$ as a}%
between normed linear spaces is itself a normed linear space under the operator norm defined above and the usual
pointwise operations of addition and scalar multiplication.  Of special importance is the family $V^* = \ofml B(V,\K)$
of the \emph{continuous} members of~$V^{\#}$, the
 \index{continuous!linear functionals}%
 \index{linear!functional!continuous}%
 \index{functional!continuous linear}%
\emph{continuous linear functionals}, also known as
 \index{bounded!linear functionals}%
 \index{linear!functional!bounded}%
 \index{functional!bounded linear}%
\emph{bounded linear functionals}.
Recall that in~\ref{def_alg_dual} we defined the algebraic dual $V^{\#}$ of a vector space~$V$. Of
greater importance in our work here is $V^*$, which is called the
 \index{<unopurstar@$V^*$ (dual space of a normed linear space)}%
 \index{V@$V^*$ (dual space of a normed linear space)}%
 \index{dual!space}%
 \index{space!dual of a normed linear}%
\df{dual space} of~$V$.  To distinguish it from the algebraic dual and the order dual, it is sometimes called the
 \index{dual!norm}%
 \index{norm!dual}%
\df{norm dual} or the
 \index{dual!topological}%
 \index{topological!dual}%
\df{topological dual} of~$V$.  (Note that this distinction is unnecessary in finite dimensional spaces
since there $V^* = V^{\#}$ (see~\ref{00015033}).
\end{exam}

\begin{prop}\label{C023228a} The family $\ofml B(V,W)$ of all bounded linear maps
 \index{B@$\ofml B(V,W)$!as a Banach space}%
 \index{Banach!space!$\ofml B(V,W)$ as a}%
between normed linear spaces is complete (with respect to the metric induced by its norm)
whenever $W$ is complete.  In particular, the dual space of every normed linear space is complete.
\end{prop}

















\begin{prop}\label{0003073} Let $U$, $V$, and $W$ be normed linear spaces.  If $S \in
\ofml B(U,V)$ and $T \in \ofml B(V,W)$, then $TS \in \ofml B(U,W)$ and $\norm{TS} \le
\norm T \norm S$.
\end{prop}

\begin{exam} On any normed linear space $V$ the
 \index{operator!identity}%
 \index{identity!operator}%
\emph{identity operator}
   \[ \id V = I_V =I\colon V \sto V \colon v \mapsto v \] is bounded and $\norm{I_V} = 1$.  The
 \index{operator!zero}%
 \index{zero!operator}%
\emph{zero operator}
   \[ \vc 0_V = \vc 0\colon V \sto V\colon v \mapsto \vc 0 \]
is also bounded and $\norm{0_V} = 0$.
\end{exam}

\begin{exam} Let $T \colon \R^2 \sto \R^3\colon (x,y) \mapsto (3x, x + 2y, x - 2y)$.
Then~$\norm{T} = \sqrt{11}$.
\end{exam}

Many students just finishing a beginning calculus course feel that differentiation is a ``nicer'' operation
than integration---probably because the rules for integration seem more complicated than those for differentiation.
However, when we regard them as linear maps exactly the opposite is true.

\begin{exam} As a linear map on the space of differentiable functions on $[0,1]$ (with the uniform norm)
integration is bounded; differentiation is not (and so is continuous nowhere!).
\end{exam}


\begin{defn} Let $f\colon V_1 \sto V_2$ be a function between normed linear spaces.  For $k = 1,2$ let
$\norm{\hphantom{M}}_k$ be the norm on~$V_k$ and $d_k$ be the metric on $V_k$ induced
by~$\norm{\hphantom{M}}_k$ (see example~\ref{0001503}).  Then $f$ is an
 \index{isometry}%
\df{isometry} (or
an \df{isometric} map) if $d_2(f(x),f(y)) = d_1(x,y)$ for all $x$, $y \in V_1$.  It is
 \index{norm!preserving}%
\df{norm preserving} if $\norm{f(x)}_2 = \norm{x}_1$ for all $x \in V_1$.)
\end{defn}

Previously in this section we have been discussing the category $\cat{NLS_{\boldsymbol\infty}}$ of normed
linear spaces and continuous linear maps.  In this category, isomorphisms preserve topological structure
(in addition to the vector space structure).  What would happen if instead we wished the isomorphisms to
preserve the metric space structure of the normed linear spaces; that is, to be isometries?  It is easy to
see that this is equivalent (in the presence of linearity) to asking the isomorphisms to be norm preserving.
In this new category, denoted by
 \index{category!$\cat{NLS_{\vc 1}}$ as a}%
 \index{NLS1@$\cat{NLS_{\vc 1}}$!the category}
$\cat{NLS_{\vc 1}}$, the morphisms would be
 \index{contractive}%
\df{contractive linear maps}, that is, linear maps $T\colon V \sto W$ between normed
linear spaces such that $\norm{Tx} \le \norm x$ for all $x \in V$.  It seems appropriate
to refer to $\cat{NLS_{\boldsymbol\infty}}$ as the
 \index{topological!category}%
 \index{category!topological}%
\emph{topological category} of normed linear spaces and to $\cat{NLS_{\vc 1}}$ as the
 \index{geometric category}%
 \index{category!geometric}%
\emph{geometric category} of normed linear spaces.  Many authors use the unmodified term
``isomorphism'' for an isomorphism in the topological category
$\cat{NLS_{\boldsymbol\infty}}$ and
 \index{isometric!isomorphism}%
 \index{isomorphism!isometric}%
``isometric isomorphism'' for an isomorphism in the geometric category $\cat{NLS_{\vc 1}}$. In these notes
we will focus on the topological category.  The isometric theory of normed linear spaces although important
is somewhat more specialized.

\begin{exer} Verify the unproved assertions in the preceding paragraph.  In particular, prove that
 \begin{enumerate}
    \item[(a)] $\cat{NLS_{\boldsymbol\infty}}$ is a category.
    \item[(b)] An isomorphism in $\cat{NLS_{\boldsymbol\infty}}$ is both a vector space isomorphism and a homeomorphism.
    \item[(c)] A linear map between normed linear spaces is an isometry if and only if it is norm preserving.
    \item[(d)] $\cat{NLS_{\vc 1}}$ is a category.
    \item[(e)] An isomorphism in $\cat{NLS_{\vc 1}}$ is both an isometry and a vector space isomorphism.
    \item[(f)] The category $\cat{NLS_{\vc 1}}$ is a subcategory of $\cat{NLS_{\boldsymbol\infty}}$
in the sense that every morphism of the former belongs to the latter.
 \end{enumerate}
\end{exer}
















\section{Finite Dimensional Spaces}
In this section are listed a few standard and useful facts about finite dimensional spaces. Proofs of these results
are quite easy to find.  Consult, for example, \cite{BrownP:1970}, section 4.4; \cite{Conway:1990}, section III.3;
\cite{Kreyszig:1978}, sections 2.4--5; or \cite{Rudin:1991}, pages 16--18.

\begin{defn} A sequence $(x_n)$ in a metric space is a
 \index{Cauchy!sequence}%
 \index{sequence!Cauchy}%
\df{Cauchy sequence} if for every $\epsilon > 0$ there exists $n_0 \in \N$ such that
$d(x_m,x_n) < \epsilon$ whenever $m$, $n \ge n_0$. This condition is often abbreviated as
follows: $d(x_m, x_n) \sto 0$ as $m$, $n \sto \infty$ (or $\lim_{m,n \sto
\infty}d(x_m,x_n) = 0$).
\end{defn}

\begin{defn} A metric space $M$ is
 \index{complete!metric space}%
 \index{metric space!complete}%
 \index{space!complete metric}%
\df{complete} if every Cauchy sequence in $M$ converges. (For a review of the most elementary facts about complete
metric spaces see the first section of chapter 20 of my online notes~\cite{Erdman:2007}.)
\end{defn}

\begin{prop}\label{prop_fdnls_Kn} If $V$ is an $n$-dimensional normed linear space over $\K$, then there exists a map
 \index{finite!dimensional normed linear spaces}%
from $V$ onto $\K^n$ which is both a homeomorphism and a vector space isomorphism.
\end{prop}

\begin{cor} Every finite dimensional normed linear space is complete.
\end{cor}

\begin{cor}\label{cor_fdvs_closed} Every finite dimensional vector subspace of a normed linear space is closed.
\end{cor}

\begin{cor}\label{cor_fdvs_norms_equiv} Any two norms on a finite dimensional vector space are equivalent.
\end{cor}

\begin{defn} A subset $K$ of a topological space $X$ is
 \index{compact}%
\df{compact} if every family $\sfml U$ of open subsets of $X$ whose union contains $K$ can be reduced to a finite
subfamily of $\sfml U$ whose union contains~$K$.
\end{defn}

\begin{prop}\label{00015032} The closed unit ball in a normed linear space is compact if and only if
the space is finite dimensional.
\end{prop}

\begin{prop}\label{00015033} If $V$ and $W$ are normed linear spaces and $V$ is finite dimensional, then
every linear map $T \colon V \sto W$ is continuous.
\end{prop}

Let $T$ be a bounded linear map between (perhaps infinite dimensional) normed linear spaces. Just as in the finite
dimensional case, the kernel and range of $T$ are objects of great importance.  In both the finite and infinite dimensional
case they are vector subspaces of the spaces which contain them.  Additionally, in both cases the kernel of $T$ is closed.
(Under a continuous function the inverse image of a closed set is closed.)  There is, however, a major difference between
the two cases.  In the finite dimensional case the range of $T$ must be closed (by~\ref{cor_fdvs_closed}).  As we will see
in example~\ref{exam_ran_nonclosed} linear maps between infinite dimensional spaces need not have closed range.






















\section{Quotients of Normed Linear Spaces}

\begin{defn}\label{quo001def} Let $A$ be an object in a concrete category $\cat C$.
A surjective morphism $A \to^\pi B$ in $\cat C$ is a
 \index{quotient!map!for concrete categories}%
 \index{pi@$\pi$ (quotient map)}%
\df{quotient map} for $A$ if a function $g\colon B \sto C$ (in $\cat{SET}$) is a morphism
(in $\cat C$) whenever $g \circ \pi$ is a morphism. An object $B$ in $\cat C$ is a
 \index{quotient!object}%
\df{quotient object} for $A$ if it is the range of some quotient map for~$A$.
\end{defn}

\begin{exam} In the category $\cat{VEC}$ of vector spaces and linear maps every surjective linear
 \index{quotient!map!in $\cat{VEC}$}%
map is a quotient map.
\end{exam}

\begin{exam}\label{C021734} In the category $\cat{TOP}$ not every epimorphism is a
 \index{quotient!map!in $\cat{TOP}$}%
quotient map.
\end{exam}

\begin{proof}[\emph{Hint for proof}] Consider the identity map on the reals with different
topologies on the domain and codomain. \ns
\end{proof}

The next item, which should be familiar from linear algebra, shows how a particular quotient object can be generated by
``factoring out a subspace''.

\begin{defn}\label{quo002def} Let $M$ be a subspace of a vector space $V$.  Define an
equivalence relation $\sim$ on $V$ by
   \[ x \sim y \qquad \text{if and only if} \qquad y - x \in M. \]
For each $x \in V$ let $[x]$ be the equivalence class containing~$x$.
 \index{<bracsqr@$[x]$ (equivalence class containing~$x$)}%
 \index{<binop@$V/M$ (quotient of $V$ by $M$)}%
Let $V/M$ be the set of all equivalence classes of elements of~$V$.  For $[x]$ and $[y]$
in $V/M$ define
  \[ [x] + [y] := [x+y]\]
and for $\alpha \in \R$ and $[x] \in V/M$ define
  \[\alpha[x] := [\alpha x].\]
Under these operations $V/M$ becomes a vector space.  It is the
 \index{quotient!vector space}%
 \index{space!quotient!vector}%
\df{quotient space} of $V$ by~$M$.  The notation $V/M$ is usually read ``$V$ mod~$M$''.
The linear map
  \[ \pi\colon V \sto V/M\colon x \mapsto [x] \]
is called the
 \index{quotient!map!for vector spaces}%
 \index{pi@$\pi$ (quotient map)}%
\df{quotient map}.
\end{defn}

\begin{exer} Verify the assertions made in definition~\ref{quo002def}. In particular,
show that $\sim$ is an equivalence relation, that addition and scalar multiplication of
the set of equivalence classes are well defined, that under these operations $V/M$ is a
vector space, that the function called here the ``quotient map'' is in fact a quotient map
in the sense of~\ref{quo001def}, and that this quotient map is linear.
\end{exer}

The following result is called the \emph{fundamental quotient theorem} or the
 \index{quotient!theorem!for $\cat{VEC}$}%
 \index{first isomorphism theorem}%
{first isomorphism theorem} for vector spaces.

\begin{thm}\label{quo005} Let $V$ and $W$ be vector spaces and $M \preceq V$. If
$T \in \ofml L(V,W)$ and $\ker T \supseteq M$, then there exists a unique $\wt T \in
\ofml L(V/M\,,W)$ which makes the following diagram commute.
 \[\xymatrix@+30pt{V \ar[d]_*+{\pi} \ar[dr]^*+{T} & \\
            V/M \ar[r]_*+{\widetilde T} & W  }\]
Furthermore, $\wt T$ is injective if and only if $\ker T = M$; and $\wt T$ is surjective
if and only if $T$ is.
\end{thm}

\begin{cor} If $T\colon V \sto W$ is a linear map between vector spaces, then
$\ran T \cong V/\ker T$.
\end{cor}

\begin{prop}\label{C023711} Let $V$ be a normed linear space and $M$ be a closed subspace
of~$V$.
 \index{quotient!normed linear space}%
 \index{normed!linear space!quotient}%
 \index{space!quotient!normed linear}%
Then the map
  \[ \norm{\hphantom{x}} \colon V/M \sto \R \colon [x]
                        \mapsto \inf\{\norm u \colon u \sim x\} \]
is a norm on~$V/M$.  Furthermore,
 \index{pi@$\pi$ (quotient map)}%
 \index{quotient!map!for normed linear spaces}%
the quotient map
  \[ \pi \colon V \sto V/M \colon x \mapsto [x] \]
is a bounded linear surjection with $\norm\pi \le 1$.
\end{prop}

\begin{defn} The norm defined on $V/M$ is called the
 \index{quotient!norm}%
 \index{norm!quotient}%
\df{quotient norm} on~$V/M$.  Under this norm $V/M$ is a normed linear space and is called the
 \index{quotient space}%
 \index{space!quotient}%
\df{quotient space} of $V$ by~$M$.  It is often referred to as ``$V$ mod $M$''.
\end{defn}

\begin{thm}[Fundamental quotient theorem for $\cat{NLS_{\boldsymbol\infty}}$]\label{C023717}
 \index{quotient!theorem!for $\cat{NLS_{\boldsymbol\infty}}$}%
Let $V$ and $W$ be normed linear spaces and $M$ be a closed subspace of~$V$.  If $T$ is a
bounded linear map from $V$ to $W$ and $\ker T \supseteq M$, then there exists a unique
bounded linear map $\wt T\colon V/M \sto W$ which makes the following diagram commute.
  \[ \xymatrix@+30pt{V \ar[d]_*+{\pi} \ar[dr]^*+{T} & \\
             V/M \ar[r]_*+{\wt T} & W  } \]
Furthermore: $\norm{\wt T} = \norm{T}$; $\wt T$ is injective if and only if $\ker T = M$;
and $\wt T$ is surjective if and only if $T$ is.
\end{thm}










































\section{Products of Normed Linear Spaces}\label{C0155}
Products and coproducts, like quotients, are best described in terms of what they \emph{do}.

\begin{defn}\label{C015511} Let $A_1$ and $A_2$ be objects in a category~$\cat C$. We say that
the triple $(P,\pi_1,\pi_2)$, where $P$ is an object and $\pi_k\colon P \sto A_k$ ($k =
1,2$) are morphisms, is a
 \index{product!in a category}%
\df{product} of $A_1$ and $A_2$ if for every object $B$ and every pair of morphisms
$f_k\colon B \sto A_k$ ($k = 1,2$) there exists a unique morphism $f\colon B \sto P$ such
that $f_k = \pi_k \circ f$ for $k = 1,2$.

It is conventional to say, ``Let $P$ be a product of \dots'' for, ``Let $(P,\pi_1,\pi_2)$
be a product of \dots''. The product of $A_1$ and $A_2$ is often written as $A_1 \times
A_2$ or as $\prod_{k = 1,2} A_k$.
\end{defn}

In a particular category products may or may not exist.  It is an interesting and
elementary fact that whenever they exist they are unique (up to isomorphism), so that we
may unambiguously speak of \emph{the} product of two objects.  When we say that a
categorical object satisfying some condition(s) is \emph{unique up to isomorphism} we
mean, of course, that any two objects satisfying the condition(s) must be isomorphic.  We
will often use the phrase
 \index{essentially!unique}%
 \index{conventions!essentially unique means unique up to isomorphism}%
 \index{uniqueness!essential}%
``essentially unique'' for ``unique up to isomorphism.''

\begin{prop} In any category products (if they exist) are
 \index{product!uniqueness of}%
 \index{uniqueness!of products}%
essentially unique.
\end{prop}

\begin{exam}\label{C015517} In the category $\cat{SET}$ the
 \index{product!in $\cat{SET}$}%
 \index{SET@$\cat{SET}$!products in}%
product of two sets $A_1$ and $A_2$ exists and is in fact the usual Cartesian product
$A_1 \times A_2$ together with the usual coordinate projections $\pi_k\colon A_1 \times
A_2 \sto A_k \colon (a_1,a_2) \mapsto a_k$.
\end{exam}

\begin{exam}\label{C015521} If $V_1$ and $V_2$ are vector spaces we make the Cartesian product
$V_1 \times V_2$ into a vector space as follows. Define addition by
  \[ (u,v) + (w,x) := (u + w, v + x)  \]
and scalar multiplication by
  \[ \alpha (u,v) := (\alpha u, \alpha v) \,. \]
This makes $V_1 \times V_2$ into a vector space and that this space together with the
usual coordinate projections $\pi_k\colon V_1 \times V_2 \sto V_k \colon (v_1,v_2)
\mapsto v_k$ ($k = 1,2$) is a
 \index{product!in $\cat{VEC}$}%
 \index{VEC@$\cat{VEC}$!products in}%
 \index{vector!space!product}%
product in the category $\cat{VEC}$. It is usually called the
 \index{direct!sum!of vector spaces}%
 \index{sum!direct}%
 \index{vector!space!direct sum}%
 \index{<binopdirectsum@$\oplus$ (direct sum)}%
\df{direct sum} of $V$ and $W$ and is denoted by~$V \oplus W$.
\end{exam}

It often possible and desirable to take the product of an arbitrary family of objects in
a category.  Following is a generalization of definition~\ref{C015511}.

\begin{defn}\label{C015524} Let $\bigl(A_\lambda\bigr)_{\lambda \in \Lambda}$ be an
indexed family of objects in a category~$\cat C$. We say that the object $P$ together
with an indexed family $\bigl(\pi_\lambda\bigr)_{\lambda \in \Lambda}$ of morphisms
$\pi_\lambda\colon P \sto A_\lambda$ is a
 \index{product!in a category}%
\df{product} of the objects $A_\lambda$ if for every object $B$ and every indexed family
$\bigl(f_\lambda\bigr)_{\lambda \in \Lambda}$ of morphisms $f_\lambda\colon B \sto
A_\lambda$ there exists a unique map $g \colon B \sto P$ such that $f_\lambda =
\pi_\lambda \circ g$ for every $\lambda \in \Lambda$.
\end{defn}

A category in which arbitrary products exist is said to be
 \index{product!complete}%
 \index{complete!product}%
\df{product complete}.  Many of the categories we encounter in these notes are product
complete.

\begin{defn}\label{C015531} Let $\bigl(S_\lambda\bigr)_{\lambda \in \Lambda}$ be an indexed
family of sets. The
 \index{Cartesian product}%
 \index{product!Cartesian}%
\df{Cartesian product} of the indexed family, denoted by
 \index{<<@$\prod S_\lambda$ (Cartesian product)}%
$\prod_{\lambda \in \Lambda} S_\lambda$ or just $\prod S_\lambda$, is the set of all
functions $f \colon \Lambda \sto \bigcup S_\lambda$ such that $f(\lambda) \in S_\lambda$
for each $\lambda \in \Lambda$.  The maps $\pi_\lambda \colon \prod S_\lambda \sto
S_\lambda \colon f \mapsto f(\lambda)$ are the canonical
 \index{coordinate!projections}%
 \index{pi@$\pi_\lambda$ (coordinate projections)}%
 \index{projection!onto coordinates}%
\df{coordinate projections}. In many cases the notation $f_\lambda$ is more convenient
than~$f(\lambda)$.  (See, for example, example~\ref{C015537} below.)
\end{defn}

\begin{exam}\label{C015534} A very important special case of the preceding definition occurs
when all of the sets $S_\lambda$ are identical: say $S_\lambda = A$ for every $\lambda
\in \Lambda$.  In this case the Cartesian product comprises all the functions which map
$\Lambda$ into~$A$. That is, $\prod_{\lambda \in \Lambda} S_\lambda = A^\Lambda$. Notice
also that in this case the coordinate projections are
 \index{evaluation!map}%
 \index{function!evaluation at a point}%
\emph{evaluation maps}. For each $\lambda$ the coordinate projection $\pi_\lambda$ takes
each point $f$ in the product (that is, each function $f$ from $\Lambda$ into~$A$) to
$f(\lambda)$ its value at~$\lambda$.  Briefly, each coordinate projection is an
evaluation map at some point.
\end{exam}

\begin{exam}\label{C015537} What is $\R^n$? It is just the set of all $n$-tuples of real numbers.
That is, it is the set of functions from $\N_n = \{1,2,\dots,n\}$ into~$\R$.  In other
words $\R^n$ is just shorthand for~$\R^{\N_n}$.  In $\R^n$ one usually writes $x_j$ for
the $j^{\text{th}}$ coordinate of a vector $x$ rather than~$x(j)$.
\end{exam}

\begin{exam} In the category $\cat{SET}$ the
 \index{product!in $\cat{SET}$}%
 \index{SET@$\cat{SET}$!products in}%
product of an indexed family of sets exists and is in fact the Cartesian product of these
sets together with the canonical coordinate projections.
\end{exam}

\begin{exam}\label{C015544} If $\bigl(V_\lambda\bigr)_{\lambda \in \Lambda}$ is an indexed
family of vector spaces we make the Cartesian product $\prod V_\lambda$ into a vector
space as follows. Define addition and scalar multiplication pointwise: for $f$, $g \in
\prod A_\lambda$ and $\alpha \in \R$
  \[ (f + g)(\lambda) := f(\lambda) + g(\lambda)  \]
and
  \[ (\alpha f)(\lambda) := \alpha f(\lambda) \,. \]
This makes $\prod V_\lambda$ into a vector space, which is sometimes called the
 \index{direct!product}%
 \index{product!direct}%
\df{direct product} of the spaces $V_\lambda$, and this space together with the canonical
coordinate projections (which are certainly  linear maps) is a
 \index{product!in $\cat{VEC}$}%
 \index{VEC@$\cat{VEC}$!products in}%
 \index{vector!space!product}%
product in the category $\cat{VEC}$.
\end{exam}

\begin{exam}\label{C023411} Let $V$ and $W$ be normed linear spaces. On the Cartesian product
$V \times W$ define
  \index{<unopn@$\norm{\hphantom{x}}_2$ (Euclidean norm or $2$-norm)}%
  \index{<unopn@$\norm{\hphantom{x}}_1$ ($1$-norm)}%
  \index{<unopn@$\norm{\hphantom{x}}_u$ (uniform norm)}%
 \begin{align*}
    \norm{(x,y)}_2 &= \sqrt{\norm x^2 + \norm y^2},  \\
    \norm{(x,y)}_1 &= \norm x + \norm y,  \\
   \intertext{and}
    \norm{(x,y)}_u &= \max\{\norm x, \norm y\}\,.
 \end{align*}
The first of these is the
 \index{Euclidean!norm!on the product of two normed spaces}%
 \index{norm!Euclidean!on the product of normed spaces}%
\df{Euclidean norm} (or
 \index{two@$2$-norm}%
 \index{norm!$2$- (on the product of normed spaces)}%
\df{$2$-norm}) on $V \times W$, the second is the
 \index{one@$1$-norm!on the product of normed spaces}%
 \index{norm!$1$- (on the product of normed spaces)}%
\df{$1$-norm}, and the last is the
 \index{uniform!norm!on the product of normed spaces}%
 \index{norm!uniform!on the product of normed spaces}%
\df{uniform norm}. Verifying that these are all norms on $V \times W$ is quite similar to
the arguments required in examples~\ref{C023125}, \ref{C023127}, and~\ref{C023129}. That
they are equivalent norms is a consequence of the following inequalities and
proposition~\ref{C023184}.
   \[ \norm x + \norm y \le \sqrt 2 \sqrt{\norm x^2 + \norm y^2} \le 2 \max\{\norm x,\norm y\} \le 2(\norm x + \norm y) \]
\end{exam}

\begin{conv}\label{C023414} In the preceding example we defined three (equivalent) norms on
the product space $V \times W$.  We will take the first of these $\norm{\;\;}_1$ as the
 \index{conventions!on choice of product norm}%
 \index{product!norm}%
 \index{norm!product}%
\df{product norm} on $V \times W$.  Thus whenever $V$ and $W$ are normed linear spaces,
unless the contrary is specified we will regard the product $V \times W$ as a normed
linear space under this norm.  Usually we write just $\norm{(x,y)}$ instead of
$\norm{(x,y)}_1$.  The product of the normed linear spaces $V$ and $W$ is usually denoted
by $V \oplus W$ and is called the
 \index{product!of normed linear spaces}%
 \index{normed!linear space!product}%
 \index{<binopdirectsum@$\oplus$ (direct sum)}%
 \index{direct!sum!of normed linear spaces}%
 \index{normed!linear space!direct sum}%
\df{direct sum} of $V$ and~$W$. (In the special case where $V = W = \R$, what metric does
the product norm induce on~$\R^2$?)
\end{conv}

\begin{exam}\label{exam_prod_nls} Show that product $V \oplus W$ of normed linear spaces is in fact a product
 \index{product!in $\cat{NLS_{\boldsymbol\infty}}$}%
 \index{NLSinf@$\cat{NLS_{\boldsymbol\infty}}$!products in}%
in the category $\cat{NLS_{\boldsymbol\infty}}$  of normed linear spaces and bounded
linear maps.
\end{exam}

\begin{prop} Let $\bigl((x_n,y_n)\bigr)$ be a sequence in the direct sum $V \oplus W$ of two normed linear
spaces. Prove that $\bigl((x_n,y_n)\bigr)$ converges to a point $(a,b)$ in $V \oplus W$ if and only
if $x_n \sto a$ in $V$ and $y_n \sto b$ in~$W$.
\end{prop}

\begin{prop}\label{C023427} Addition is a continuous operation on a normed linear space~$V$.
 \index{continuity!of addition}%
 \index{addition!continuity of}%
That is, the map
   \[ A \colon V \oplus V \sto V \colon (x,y) \mapsto x + y \]
is continuous.
\end{prop}

\begin{exer} Give a \emph{very} short proof (no $\epsilon$'s or $\delta$'s or open sets) that if
$(x_n)$ and $(y_n)$ are sequences in a normed linear space which converge to $a$ and $b$,
respectively, then $x_n + y_n \sto a + b$.
\end{exer}

\begin{prop}\label{C023434} Scalar multiplication is a continuous operation on a normed
 \index{continuity!of scalar multiplication}%
 \index{scalar!multiplication!continuity of}%
linear space~$V$ in the sense that the map
   \[ S \colon \K \times V \sto V \colon  (\alpha,x) \mapsto \alpha x \]
is continuous.
\end{prop}

\begin{prop} If $B$ and $C$ are subsets of a normed linear space and $\alpha$ is a scalar, then
 \begin{enumerate}
   \item $\clo{\alpha B} = \alpha \clo B$; and
   \item $\clo B + \clo C  \subseteq  \clo{B + C}$.
 \end{enumerate}
\end{prop}

\begin{exam} If $B$ and $C$ are closed subsets of a normed linear space, then it does not
necessarily follow that $B + C$ is closed (and therefore $\clo B + \clo C$ and $\clo{B +
C}$ need not be equal).
\end{exam}

\begin{prop} Let $C_1$ and $C_2$ be compact subsets of a normed linear space~$V$. Then
   \begin{enumerate}
     \item $C_1 \times C_2$ is a compact subset of $V \times V$, and
     \item $C_1 + C_2$ is a compact subset of~$V$.
   \end{enumerate}
\end{prop}

\begin{exam}\label{C023441}  Let $X$ be a topological space.  Then the family
$\fml C(X)$ of all continuous complex valued functions on~$X$ is a
 \index{C@$\fml C(X)$!as a vector space}%
vector space under the usual pointwise operations of addition and scalar multiplication.
\end{exam}

\begin{exam}\label{C023443}  Let $X$ be a topological space.  We denote
 \index{C@$\fml C_b(X)$!continuous bounded functions on $X$}%
by $\fml C_b(X)$ the family of all bounded continuous complex valued functions on~$X$.  It is a
 \index{C@$\fml C_b(X)$!as a normed linear space}%
 \index{normed!linear space!$\fml C_b(X)$ as a}%
normed linear space under the uniform norm.  In fact, it is a subspace of~$\fml B(X)$
(see example~\ref{C023131}).
\end{exam}

\begin{exam}\label{C023464} If $S$ is a set and $a \in S$, then the evaluation functional
     \[ E_a\colon \fml B(S) \sto \R\colon f \mapsto f(a) \]
is a bounded linear functional on the space $\fml B(S)$ of all bounded real valued
functions on~$S$. If $S$ happens to be a topological space then we may also regard $E_a$
as a member of the dual space of $\fml C_b(S)$.  (In either case, what is~$\norm{E_a}$?)
\end{exam}

\begin{defn}\label{def_norm_alg} Let $A$ be an algebra on which a norm has been defined.  Suppose that additionally the
 \index{submultiplicative}%
\df{submultiplicative} property
  \[ \norm{xy} \le \norm x\,\norm y \]
is satisfied for all $x$, $y \in A$.  Then $A$ is a
 \index{normed!algebra}%
 \index{algebra!normed}%
\df{normed algebra}.  We make one further requirement: if the algebra $A$ is unital, then
   \[ \norm{\vc 1} = 1\,. \]
In this case $A$ is a
 \index{unital!normed algebra}%
 \index{normed!algebra!unital}%
 \index{algebra!unital normed}%
\df{unital normed algebra}.
\end{defn}

\begin{exam}\label{C023448} In example~\ref{C023131} we showed that the family
 \index{B@$\fml B(S)$!as a normed algebra}%
 \index{normed!algebra!$\fml B(S)$ as a}%
$\fml B(S)$ of bounded real valued functions on a set $S$ is a normed linear space under
the uniform norm.  It is also a commutative unital normed algebra.
\end{exam}

\begin{prop} Multiplication is a continuous operation on a normed algebra~$V$ in the sense
that the map
   \[ M \colon V \times V \sto V \colon (x,y) \mapsto  xy \]
is continuous.
\end{prop}

\begin{proof}[\emph{Hint for proof}] If you can prove that multiplication on the real numbers
is continuous you can almost certainly prove that it is continuous on arbitrary normed
algebras. \ns
\end{proof}

\begin{exam}\label{C023454} The family $\fml C_b(X)$ of bounded continuous real valued
 \index{C@$\fml C_b(X)$!as a normed algebra}%
 \index{normed!algebra!$\fml C_b(X)$ as a}%
functions on a topological space~$X$ is, under the usual pointwise operations and the
uniform norm, a commutative unital normed algebra. (See example~\ref{C023443}.)
\end{exam}

\begin{exam}\label{C023456} The family $\ofml B(V)$ of all (bounded linear) operators
 \index{B@$\ofml B(V)$!as a unital normed algebra}%
 \index{normed!algebra!$\ofml B(V)$ as a}%
on a normed linear space $V$ is a unital normed algebra. (See~\ref{C023228} and~\ref{C023228a}.)
\end{exam}





























\section{Coproducts of Normed Linear Spaces}\label{C0156}
Coproducts are like products---except all the arrows are reversed.

\begin{defn}\label{C015611} Let $A_1$ and $A_2$ be objects in a category~$\cat C$.  The triple
$(Q,\iota_1,\iota_2)$, ($Q$ is an object and $\iota_k\colon A_k \sto Q$, $k = 1,2$ are morphisms) is a
 \index{coproduct}%
\df{coproduct} of $A_1$ and $A_2$ if for every object $B$ and every pair of morphisms
$f_k\colon A_k \sto B$ (k = 1,2) there exists a unique morphism $f\colon P \sto B$ such
that $f_k = f \circ \iota_k$ for $k = 1,2$.
\end{defn}

\begin{prop} In any category coproducts (if they exist) are
 \index{coproduct!uniqueness of}%
essentially unique.
\end{prop}

\begin{defn}\label{C0006127} Let $A$ and $B$ be sets.  When $A$ and $B$ are disjoint we will
often use the notation
 \index{<binopuniond@$\uplus$, $\biguplus$ (disjoint union)}%
$A \uplus B$ instead of $A \cup B$ (to emphasize the disjointness of $A$ and~$B$). When
$C = A \uplus B$ we say that $C$ is the
 \index{disjoint!union}%
 \index{union!disjoint}%
\df{disjoint union} of $A$ and~$B$.  Similarly we frequently choose to write the union of
a \emph{pairwise disjoint} family $\sfml A$ of sets as $\biguplus\sfml A$. And the
notation $\biguplus_{\lambda \in \Lambda} A_\lambda$ may be used to denote the union of a
\emph{pairwise disjoint} indexed family $\{A_\lambda\colon  \lambda \in \Lambda\}$ of
sets. When $C = \biguplus\sfml A$ or $C = \biguplus_{\lambda \in \Lambda} A_\lambda$ we
say that $C$ is the \df{disjoint union} of the appropriate family.
\end{defn}

\begin{defn} In the preceding definition we introduced the notation $A \uplus B$ for the
union of two disjoint sets $A$ and~$B$.  We now extend this notation somewhat and allow
ourselves to take the \emph{disjoint union} of sets $A$ and $B$ even if they are not
disjoint. We ``make them disjoint'' by identifying (in the obvious way) the set $A$ with the
set $A' = \{(a,1)\colon a \in A\}$ and $B$ with the set $B' = \{(b,2)\colon b \in B\}$.  Then we
denote the union of $A'$ and $B'$, which \emph{are} disjoint, by $A \uplus B$ and call it the
 \index{disjoint!union}%
 \index{union!disjoint}%
 \index{<binopuniond@$\uplus$, $\biguplus$ (disjoint union)}%
\df{disjoint union} of $A$ and~$B$.

In general, if $\bigl(A_\lambda\bigr)_{\lambda \in \Lambda}$ is an indexed family of
sets, its
 \index{disjoint!union}%
 \index{union!disjoint}%
\df{disjoint union}, \smash[b]{$\biguplus\limits_{\lambda \in \Lambda}A_\lambda$} is
defined to be $\bigcup \{(A_\lambda,\lambda)\colon \lambda \in \Lambda\}$.
\end{defn}

\begin{exam}\label{C015617} In the category $\cat{SET}$ the
 \index{coproduct!in $\cat{SET}$}%
 \index{SET@$\cat{SET}$!coproducts in}%
coproduct of two sets $A_1$ and $A_2$ exists and is their disjoint union $A_1 \uplus A_2$
together with the obvious inclusion maps $\iota_k\colon A_k  \sto A_1 \uplus A_2$ ($k = 1,2$).
\end{exam}

It is interesting to observe that while the product and coproduct of a finite collection
of objects in the category $\cat{SET}$ are quite different, in the more complex category
$\cat{VEC}$ they turn out to be exactly the same thing.

\begin{exam} If $V_1$ and $V_2$ are vector spaces make the Cartesian product $V_1 \times V_2$
into a vector space as in example~\ref{C015521}. This space together with the obvious
injections is a
 \index{coproduct!in $\cat{VEC}$}%
 \index{VEC@$\cat{VEC}$!coproducts in}%
 \index{vector!space!coproduct}%
coproduct in the category $\cat{VEC}$.
\end{exam}

It often possible and desirable to take the coproduct of an arbitrary family of objects
in a category.  Following is a generalization of definition~\ref{C015611}.

\begin{defn}\label{C015624} Let $\bigl(A_\lambda\bigr)_{\lambda \in \Lambda}$ be an indexed
family of objects in a category~$\cat C$. We say that the object $C$ together with an
indexed family $\bigl(\iota_\lambda\bigr)_{\lambda \in \Lambda}$ of morphisms
$\iota_\lambda\colon A_\lambda \sto C$ is a
 \index{coproduct!in a category}%
\df{coproduct} of the objects $A_\lambda$ if for every object $B$ and every indexed
family $\bigl(f_\lambda\bigr)_{\lambda \in \Lambda}$ of morphisms $f_\lambda\colon
A_\lambda \sto B$ there exists a unique map $g \colon C \sto B$ such that $f_\lambda = g
\circ \iota_\lambda$ for every $\lambda \in \Lambda$.  The usual notation for the
coproduct of the objects $A_\lambda$ is
 \index{<<@$\coprod A_\lambda$ (coproduct)}%
$\coprod_{\lambda \in \Lambda} A_\lambda$.
\end{defn}

\begin{exam} In the category $\cat{SET}$ the
 \index{coproduct!in $\cat{SET}$}%
 \index{SET@$\cat{SET}$!coproducts in}%
coproduct of an indexed family of sets exists and is the disjoint union of these sets.
\end{exam}

\begin{defn}  Let $S$ be a set and $V$ be a vector space. The
 \index{support!of a function}%
\df{support} of a function $f \colon S \sto V$ is $\{s \in S \colon f(s) \ne \vc 0\}$.
\end{defn}

\begin{exam}\label{X_dir_sum_VEC} If $\bigl(V_\lambda\bigr)_{\lambda \in \Lambda}$ is an indexed family of
vector spaces we make the Cartesian product into a vector space as in example~\ref{C015544}. The set of functions $f$
belonging to $\prod V_\lambda$ which have finite support (that is, which are nonzero only finitely often) is clearly
a subspace of $\prod V_\lambda$.  This subspace is the
 \index{direct!sum}%
 \index{sum!direct}%
\df{direct sum} of the spaces $V_\lambda$.  It is denoted by $\bigoplus_{\lambda \in
\Lambda} V_\lambda$.  This space together with the obvious injective maps (which are
linear) is a
 \index{coproduct!in $\cat{VEC}$}%
 \index{VEC@$\cat{VEC}$!coproducts in}%
 \index{vector!space!coproduct}%
coproduct in the category $\cat{VEC}$.
\end{exam}

\begin{exer}  What are the coproducts in the category $\cat{NLS_{\boldsymbol\infty}}$ ?
\end{exer}

\begin{exer} Identify the product and coproduct of two spaces $V$ and $W$ in the category
$\cat{NLS_{\vc 1}}$ of normed linear spaces and linear contractions.
\end{exer}













\section{Nets}\label{C0314}
Nets (sometimes called \emph{generalized sequences}) are useful for characterizing many
properties of general topological spaces.  In such spaces they are used in basically
the same way as sequences are used in metric spaces.

\begin{defn} A preordered set in which every pair of elements has an upper bound is a
 \index{directed!set}%
 \index{set!directed}%
\df{directed set}.
\end{defn}

\begin{exam}\label{C031414}  If $S$ is a set, then $\fin S$ is a directed set under
inclusion $\subseteq$.
\end{exam}

\begin{defn} Let $S$ be a set and $\Lambda$ be a directed set.  A mapping
$x\colon \Lambda \sto S$ is a
 \index{net}%
\df{net} in $S$ (or a \emph{net of members of $S$}). The value of $x$ at $\lambda \in
\Lambda$ is usually written $\sbsb x\lambda$ (rather than $x(\lambda)$) and the net $x$
itself is often denoted by $\bigl(\sbsb x\lambda\bigr)_{\lambda \in \Lambda}$ or just
$\bigl(\sbsb x\lambda\bigr)$.
\end{defn}

\begin{exam} The most obvious examples of nets are sequences.  These are functions whose domain is the (preordered) set $\N$ of natural numbers.
\end{exam}

\begin{defn} A net $x = \bigl(\sbsb x\lambda\bigr)$ is said to
 \index{eventually}%
\df{eventually} have some specified property if there exists $\lambda_0 \in \Lambda$ such that $\sbsb x\lambda$ has that property
whenever $\lambda \ge \lambda_0$; and it has the property
 \index{frequently}%
\df{frequently} if for every $\lambda_0 \in \Lambda$ there exists $\lambda \in \Lambda$ such that $\lambda \ge \lambda_0$
and $\sbsb x\lambda$ has the property.
\end{defn}

\begin{defn} A
 \index{neighborhood}%
\df{neighborhood} of a point in a topological space is any set which contains an open set containing the point.
\end{defn}

\begin{defn}  A net $x$ in a topological space $X$
 \index{convergence!of nets}%
 \index{net!convergence of a}%
\df{converges} to a point $a \in X$ if it is eventually in every neighborhood of~$a$. In
this case $a$ is the
 \index{limit!of a net}%
 \index{net!limit of a}%
\df{limit} of $x$ and we write
   \[ x_{{}_{\sst \lambda}} \to_{\lambda \in \Lambda}a \]
or just
 \index{<arrowto@$x_\lambda \sto a$ (convergence of nets)}%
$x_{{}_{\sst \lambda}} \sto a$.  When limits are unique we may also use the notation
$\lim_{\lambda \in \Lambda} \sbsb x\lambda = a$ or, more simply, $\lim \sbsb x\lambda =
a$.

The point $a$ in $X$ is a
 \index{cluster point!of a net}%
 \index{net!cluster point of a}%
\df{cluster point} of the net $x$ if $x$ is frequently in every neighborhood of~$a$.
\end{defn}

\begin{defn} Let $(X,\sfml T)$ be a topological space.  A subfamily $\sfml S \subseteq
\sfml T$ is a
 \index{subbase}%
\df{subbase} for the topology $\sfml T$ if the family of all finite intersections of
members of $\sfml S$ is a base for~$\sfml T$.
\end{defn}

\begin{exam}\label{C031429} Let $X$ be a topological space and $a \in X$. A
 \index{neighborhood!base for}%
 \index{base!for neighborhoods of a point}%
\df{base} for the family of neighborhoods of the point $a$ is a family $\sfml B_a$ of
neighborhoods of $a$ with the property that every neighborhood of $a$ contains at least
one member of~$\sfml B_a$.  In general, there are many choices for a neighborhood base at
a point. Once such a base is chosen we refer to its members as
 \index{basic!neighborhoods}%
 \index{neighborhood!basic}%
\emph{basic neighborhoods} of~$a$.

Let $\sfml B_a$ be a base for the family of neighborhoods at $a$.  Order $\sfml B_a$ by
containment (the reverse of inclusion); that is, for $U$, $V \in \sfml B_a$ set
  \[ U \preceq V \text{ if and only if } U \supseteq V\,. \]
This makes $\sfml B_a$ into a directed set. Now (using the axiom of choice) choose one
element $x_{{}_U}$ from each set $U \in \sfml B_a$.  Then $(\sbsb xU)$ is a net in $X$
and $\sbsb xU \sto a$.
\end{exam}

The next two propositions assure us that in order for a net to converge to a point $a$ in
a topological space it is sufficient that it be eventually in every basic (or even
subbasic) neighborhood of~$a$.

\begin{prop}\label{C031431} Let $\sfml B$ be a base for the topology on a topological
space~$X$ and $a$ be a point of~$X$. A net $(x_\lambda)$  converges to  $a$ if and only
if it is eventually in every neighborhood of~$a$ which belongs to~$\sfml B$.
\end{prop}

\begin{prop}\label{C031434} Let $\sfml S$ be a subbase for the topology on a topological
space~$X$ and $a$ be a point of~$X$.  A net $(x_\lambda)$ converges to $a$ if and only if
it is eventually in every neighborhood of~$a$ which belongs to~$\sfml S$.
\end{prop}

\begin{exam}\label{C031437} Let $J = [a,b]$ be a fixed interval in the real line,
$\vc x = (x_0,x_1,\dots,x_n)$ be an $(n+1)$-tuple of points of $J$, and $\vc t =
(t_1,\dots,t_n)$ be an $n$-tuple. The pair $(\vc x;\vc t)$ is a
 \index{partition!with selection}%
 \index{selection}%
\df{partition with selection} of the interval $J$ if:
 \begin{enumerate}
  \item $x_{k-1} < x_k$ for $1 \le k \le n$;
  \item $t_k \in [x_{k-1},x_k]$ for $1 \le k \le n$;
  \item $x_0 = a$; and
  \item $x_n = b$.
 \end{enumerate}
The idea is that $\vc x$ \emph{partitions} the interval into subintervals and $t_k$ is
the point \emph{selected} from the $k^{\text{th}}$ subinterval.  If $\vc x =
(x_0,x_1,\dots,x_n)$, then $\{\vc x\}$ denotes the \emph{set} $\{x_0,x_1,\dots,x_n\}$.
Let $P = (\vc x;\vc s)$ and $Q = (\vc y;\vc t)$ be partitions (with selections) of the
interval~$J$. We write $P \preccurlyeq Q$ and say that $Q$ is a
 \index{refinement}%
\df{refinement} of $P$ if $\{\vc x\} \subseteq \{\vc y\}$. Under the relation
$\preccurlyeq$ the family $\sfml P$ of partitions with selection on $J$ is a directed
(but \emph{not} partially ordered) set.

Now suppose $f$ is a bounded function on the interval $J$ and $P = (\vc x;\vc t)$ is a
partition of $J$ into $n$ subintervals. Let $\Delta x_k := x_k - x_{k-1}$ for $1 \le k
\le n$. Then define
   \[ S_f(P) := \sum_{k=1}^n f(t_k)\,\Delta x_k\,. \]
Each such sum $S_f(P)$ is a
 \index{Riemann!sum}%
 \index{sum!Riemann}%
\df{Riemann sum} of $f$ on~$J$. Notice that since $\sfml P$ is a directed set, $S_f$ is a
net of real numbers.  If the net $S_f$ of Riemann sums converges we say that the function
$f$ is
 \index{Riemann!integrable}%
 \index{integrable!Riemann}%
\df{Riemann integrable}.  The limit of the net $S_f$ (when it exists) is the
 \index{Riemann!integral}%
 \index{integral!Riemann}%
\df{Riemann integral} of $f$ over the interval~$J$ and is denoted by $\int_a^b f(x)\,dx$.
\end{exam}

\begin{exam}\label{exam_net_indiscrete} In a topological space $X$ with the
indiscrete topology (where the only open sets are $X$ and the empty set) every net in the space
converges to every point of the space.
\end{exam}

\begin{exam} In a topological space $X$ with the discrete topology (where every subset of $X$ is open) a net
in the space  converges if and only if it is eventually constant.
 \end{exam}

\begin{defn} A
 \index{Hausdorff}%
 \index{separation!axiom!Hausdorff}%
 \index{topological!space!Hausdorff}%
 \index{topology!Hausdorff}%
\df{Hausdorff} topological space is one in which every pair of distinct points can be
separated by open sets; that is, for every pair $x$ and $y$ of distinct points in the space there exist
open neighborhoods of $x$ and $y$ that are disjoint.
\end{defn}

As example~\ref{exam_net_indiscrete} shows nets in general topological spaces need not have unique
limits. A space does not require very restrictive assumptions however to make this
pathology go away. For limits to be unique it is sufficient to require that a space be
Hausdorff. Interestingly, it turns out that this condition is also necessary.

\begin{prop}\label{C031441} If $X$ is a topological space, then the following are equivalent:
 \begin{enumerate}
  \item $X$ is Hausdorff.
  \item No net in $X$ can converge to more than one point.
  \item The diagonal in $X \times X$ is closed.
 \end{enumerate}
\end{prop}
(The
 \index{diagonal}%
\df{diagonal} of a set $S$ is $\{(s,s) \colon s \in S\} \subseteq S \times S$.)

\begin{exam} Recall from beginning calculus that the $n^{\text{th}}$
 \index{partial!sum}%
 \index{sum!partial}%
 \index{infinite series!partial sum of an}%
 \index{series!infinite}%
\df{partial sum} of an infinite series $\sum_{k=1}^\infty x_k$ is defined to be $s_n =
\sum_{k=1}^n x_k$.  The series is said to
 \index{convergence!of infinite series}%
 \index{infinite series!convergence of an}%
\df{converge} if the sequence $(s_n)$ of partial sums converges and, if the series
converges, its
 \index{sum!of an infinite series}%
 \index{infinite series!sum of an}%
 \index{series!sum of a}%
\emph{sum} is the limit of the sequence $(s_n)$ of partial sums.  Thus, for example, the
infinite series $\sum_{k=1}^\infty 2^{-k}$ converges and its sum is~$1$.
\end{exam}

The idea of \emph{summing} an infinite series depends heavily on the fact that the
partial sums of the series are linearly ordered by inclusion. The ``partial sums'' of a
net of numbers need not be partially ordered, so we need a new notion of
``summation'', sometimes referred to as \emph{unordered summation}.

\begin{notn} If $A$ is a set we denote
 \index{fin@$\fin A$ (family of finite subsets of a set $A$)}%
by $\fin A$ the family of all finite subsets of~$A$. (Note that this is a directed set under the
relation~$\subseteq$.)
\end{notn}

\begin{defn}\label{def_summable_R} Let $A \subseteq \R$.  For every $F \in \fin A$ define $S_F$ to be the sum of
the numbers in~$F$, that is, $S_F = \sum F$.  Then $S = \bigl(S_F\bigr)_{F \in \fin A}$
is a net in~$\R$.  If this net converges, the set $A$ of numbers is said to be
 \index{summable!set of real numbers}%
\df{summable}; the limit of the net is the
 \index{sum!of a summable set}%
\df{sum} of $A$ and is denoted by $\sum A$.
\end{defn}

\begin{defn}  A net $(x_\lambda)$ of real numbers is
 \index{increasing!net in $\R$}%
 \index{net!increasing}%
\df{increasing} if $\lambda \le \mu$ implies $x_\lambda \le x_\mu$.  A net of real (or complex)
numbers is
 \index{bounded!net in $\K$}%
 \index{net!bounded}%
\df{bounded} if its range is.
\end{defn}

\begin{exam}\label{C031454} Let $f(n) = 2^{-n}$ for every $n \in \N$, and let $x \colon
\fin(\N) \sto \R$ be defined by $x(A) = \sum_{n \in A}f(n)$.
 \begin{enumerate}
  \item The net $x$ is increasing.
  \item The net $x$ is bounded.
  \item The net $x$ converges to 1.
  \item The set of numbers $\{\frac12,\frac14,\frac18, \frac1{16},\dots\}$ is summable and
its sum is~$1$.
  \end{enumerate}
\end{exam}

\begin{exam} In a metric space every convergent sequence is bounded.  Give an example to show
that a convergent net in a metric space (even in~$\R$) need not be bounded.
\end{exam}

It is familiar from beginning calculus that bounded increasing sequences in $\R$
converge. The preceding example is a special case of a more general observation: bounded
increasing nets in $\R$ converge.

\begin{prop}\label{C031461} Every bounded increasing net
$\bigl(x_{{}_{\sst \lambda}}\bigr)_{\lambda \in \Lambda}$ of real numbers converges and
  \[ \lim x_{{}_{\sst \lambda}} = \sup\{x_{{}_{\sst \lambda}}\colon \lambda \in \Lambda\}\,. \]
\end{prop}

\begin{exam} Let $(x_k)$ be a sequence of distinct positive real numbers and $A = \{x_k \colon k \in \N\}$.
Then $\sum A$ (as defined above) is equal to the sum of the series $\sum_{k=1}^\infty
x_k$.
\end{exam}

\begin{exam} The set $\bigl\{\frac1k \colon k \in \N \bigr\}$ is not summable and the infinite
series $\sum_{k=1}^\infty \frac1k$ does not converge.
\end{exam}

\begin{exam} The set $\bigl\{(-1)^k \frac1k\colon k \in \N \bigr\}$ is not summable but the
infinite series $\sum_{k=1}^\infty (-1)^k \frac1k$ does converge.
\end{exam}

\begin{prop} Every summable subset of real numbers is countable.
\end{prop}

\begin{proof}[\emph{Hint for proof.}] If $A$ is a summable set of real numbers and $n \in \N$,
how many members of $A$ can be larger than~$1/n$?  \ns
\end{proof}

In proposition~\ref{C031441} we have already seen one example of a way in which nets
characterize a topological property: a space is Hausdorff if and only if limits of nets
in the space are unique (when they exist).  Here is another topological
property---continuity at a point---that can be conveniently characterized in terms of
nets.

\begin{prop}\label{C031471}  Let $X$ and $Y$ be topological spaces.  A function
 \index{continuity!characterization by nets}%
$f\colon X \sto Y$ is continuous at a point $a$ in $X$ if and only if $f\bigl(\sbsb x\lambda\bigr) \sto f(a)$ in $Y$
whenever $\bigl(\sbsb x\lambda\bigr)$ is a net converging to $a$ in~$X$.
\end{prop}

In metric spaces we characterize closures and interiors of sets in terms of sequences. For
general topological spaces we must replace sequences by nets.

In defining the terms ``interior'' and ``closure'' we have choices similar to the ones we had in defining ``convex hull'' in
\ref{smplx005def}.  There is an 'abstract' definition and there is a `constructive' one.  You may find it helpful to recall
proposition~\ref{smplx006} and the paragraph preceding it.

\begin{defn}\label{def_interior} The
 \index{interior}%
\df{interior} of a set $A$ in a topological space $X$ is the largest open subset of $X$ contained in~$A$; that is, the
union of all open subsets of $X$ contained in~$A$. The interior of $A$
 \index{<unopur@$\intr A$ (interior of~$A$)}%
is denoted by $\intr A$.
\end{defn}

\begin{prop} The preceding definition makes sense.  Furthermore, a point $a$ in a topological space $X$ is in the interior of a subset
 \index{interior!characterization of by nets}%
 \index{characterization!by nets!of interior}%
 \index{nets!characterization!of interior by}%
$A$ of $X$ if and only if every net in $X$ that converges to $a$ is eventually in~$A$.
\end{prop}

\begin{proof}[\emph{Hint for proof}]  One direction is easy. In the other use proposition~\ref{C031431} and the sort
of net constructed in example~\ref{C031429}.   \ns
\end{proof}

\begin{defn}\label{def_closure} The
 \index{closure}%
\df{closure} of a set $A$ in a topological space $X$ is the smallest closed subset of $X$ containing~$A$; that is, the
intersection of all closed subsets of $X$ which contain~$A$.  The closure of $A$
 \index{<unopu@$\clo A$ (closure of~$A$)}%
is denoted by~$\clo A$.
\end{defn}

\begin{prop}\label{C031481} The preceding definition makes sense.  Furthermore, a point $b$ in a topological space $X$ belongs to the
 \index{closure!characterization of by nets}%
 \index{charactrization!by nets!of closure}%
 \index{nets!characterization!of closure by}%
closure of a subset $A$ of~$X$ if and only if some net in $A$ converges to~$b$.
\end{prop}

\begin{cor} A subset $A$ of a topological space is closed if and only if the limit of every convergent net in $A$ belongs to~$A$.
\end{cor}


















\section{The Hahn-Banach Theorems}

It should be said at the beginning that referring to ``the'' \emph{Hahn-Banach theorem} is something of a misnomer.
A writer who claims to be using the \emph{Hahn-Banach theorem} may be using any of a dozen or so versions of the
theorem or its corollaries.  What do the various versions have in common?  Basically they assure us that normed
linear spaces have, in general, a rich enough supply of bounded linear functionals to make their dual spaces useful
tools in studying the spaces themselves. That is, they guarantee an interesting and productive duality theory.
Below we consider four versions of the `theorem' and three of its corollaries.  We will succumb to usual practice:
when in the sequel we make use of any of these results, we will say we are using `the' \emph{Hahn-Banach theorem}.

As is often the case in mathematics the attribution of these results to Hahn and Banach is a bit peculiar.  It seems
that the basic insights were due originally to Eduard Helly.  For an interesting history of the `Hahn-Banach' results
see~\cite{Hochstadt:1980}.

Incidentally, in the name ``Banach'', the ``a'''s are pronounced as in \emph{car} and the ``ch'' as in \emph{chaos};
and the accent is on the first syllable (BAH-nakh).

\begin{defn} A real valued function $f$ on a real vector space $V$ is a
 \index{sublinear}%
\df{sublinear functional} if $f(x + y) \le f(x) + f(y)$ and $f(\alpha x) = \alpha f(x)$
for all $x$, $y \in V$ and $\alpha \ge 0$.
\end{defn}

\begin{exam} A norm on a real vector space is a sublinear functional.
\end{exam}

Our first version of the \emph{Hahn-Banach theorem} is pure linear algebra and applies only to \emph{real} vector spaces.
It allows us to extend certain linear functionals on such spaces from subspaces to the whole space.

\begin{thm}[Hahn-Banach theorem I]\label{HBTI} Let $M$ be a subspace of a real vector space $V$ and $p$
 \index{Hahn-Banach theorem}%
be a sublinear functional on~$V$. If $f$ is a linear functional on $M$ such that $f \le p$ on $M$, then
$f$ has an extension $g$ to all of $V$ such that $g \le p$ on~$V$.
\end{thm}

\begin{proof}[\emph{Hint for proof}] In this theorem the notation $h \le p$ means that $\dom h \subseteq \dom p = V$
and that $h(x) \le p(x)$ for all $x \in \dom h$.  In this case we say that $h$ is
 \index{dominate}%
\df{dominated} by~$p$.  We write
 \index{<binrelle@$f \preccurlyeq g$ ($g$ is an extension of~$f$)}%
$f \preccurlyeq g$ if $g$ is an extension of~$f$ (that is,
if $f(x) = g(x)$ for all $x \in \dom f$).

Let $\fml S$ be the family of all real valued linear functionals on subspaces of $V$ which are extensions of $f$ and are
dominated by~$p$. The family $\fml S$ is partially ordered by $\preccurlyeq$.  Use \emph{Zorn's lemma} to produce a
maximal element $g$ of~$\fml S$.  The proof is finished if you can show that $\dom{g} = V$.

To this end, suppose to the contrary that there is a vector $c \in V$ which does not belong to~$D = \dom{g}$.  Prove first that
   \begin{equation}\label{eq_HBTI}
      g(y) - p(y - c) \le - g(x) + p(x + c)
   \end{equation}
for all $x$, $y \in D$.  Conclude that there is a number $\gamma$ which lies between the supremum of all the numbers
on the left side of~\eqref{eq_HBTI} and the infimum of those on the right.

Let $E$ be the span of $D \cup \{c\}$.  Define $h$ on $E$ by $h(x + \alpha c) = g(x) + \alpha\gamma$ for $x \in D$ and
$\alpha \in \R$.  Now show $h \in \fml S$.  (Take $\alpha > 0$ and work separately with $h(x + \alpha c)$ and $h(x - \alpha c)$.) \ns
\end{proof}

The next theorem is the version of theorem~\ref{HBTI} for spaces with complex scalars.  It is sometimes referred to as the
 \index{Bohnenblust-Sobczyk-Suhomlinov theorem}%
\emph{Bohnenblust-Sobczyk-Suhomlinov theorem}.

\begin{thm}[Hahn-Banach theorem II]\label{HBTII} Let $V$ be a complex vector space, $p$ be a seminorm on~$V$, and $M$ be a subspace of~$V$.
If $f\colon M \sto \C$ is a linear functional such that $\abs f \le p$, then $f$ has an extension $g$ which is a linear functional
on all of~$V$ satisfying $\abs{g} \le p$.
\end{thm}

\begin{proof}[\emph{Hint for proof}] This result is essentially a straightforward corollary to the real version of the
\emph{Hahn-Banach theorem} given above in~\ref{HBTI}.  Let $V_{\R}$ be the real vector space associated with~$V$ (same vectors
but using only $\R$ as the scalar field).  Notice that if $u$ is the real part of $f$ (that is,
 \index{R@$\Re(z)$ (the real part of $z \in \C$)}%
$u(x) = \Re(f(x))$ for each $x \in M$), then $u$ is an $\R$-linear functional on $M$ (that is, $f(x + y) = f(x) + f(y)$ and $f(\alpha x) = \alpha f(x)$
for all $x$, $y \in M$ and all $\alpha \in \R$) and $f(x) = u(x) - i\,u(ix)$.  And conversely, if $u$ is a $\R$-linear functional on some vector space,
then $f(x) := u(x) - i\,u(ix)$ defines a (complex valued) linear functional on that space.  Notice also that if $u$ is the real part of~$f$,
then $\abs{u(x)} \le p$ for all $x$ if and only if $\abs{f(x)} \le p$ for all~$x$.  With these observations in mind just apply~\ref{HBTI}
to the real part of $f$ on~$V_\R$.  \ns
\end{proof}

Next is perhaps the most often used version of the \emph{Hahn-Banach theorem}.

\begin{thm}[Hahn-Banach theorem III]\label{HBTIII} Let $M$ be a vector subspace of a (real or
 \index{Hahn-Banach theorem}%
complex) normed linear space $V$. Then every continuous linear functional $f$ on $M$ has
an extension $g$ to all of $V$ such that $\norm{g} = \norm f$.
\end{thm}

\begin{proof}[\emph{Hint for proof}] Consider the function $p\colon V \sto \R \colon x \mapsto \norm f \norm x$.  \ns
\end{proof}

It is sometimes useful, when proving a theorem, to consider the possibility of more ``natural'' or even  more ``obvious'' proofs.
On occasion one is rewarded by an enlightening simplification and clarification of complex material.  And sometimes one is fooled.

\begin{exer} Your buddy Fred R. Dimm thinks he has discovered a really simple direct proof of
the version of the Hahn-Banach theorem given in proposition~\ref{HBTIII}: A continuous
linear functional $f$ defined on a (linear) subspace $M$ of a normed linear space~$V$ can
be extended to a continuous linear functional $g$ on all of $V$ without increasing
its norm. He says that all one needs to do is let $N$ be any vector space complement
of~$M$ and let $B$ be a (Hamel) basis for $N$.  Define $g(e) = 0$ for every vector $e
\in B$ and $g = f$ on~$M$. Then extend $g$ by linearity to all of~$V$. Surely,
Fred says, $g$ is linear by construction and its norm has not increased since we
have made it zero where it was not previously defined.  Show Fred the error of his ways.
\end{exer}

\begin{thm}[Hahn-Banach theorem IV]\label{HBTIV} Let $V$ be a (real or complex) normed
 \index{Hahn-Banach theorem}%
linear space, $M$ be a vector subspace of $V$, and suppose that $y \in V$ is such that
 \index{dist@$\dist(y,M)$ (the distance between $y$ and $M$)}%
$d := \dist(y,M) > 0$.  Then there exists a functional $g \in V^*$ such that $g^{\sto}(M) = \{0\}$,
$g(y) = d$, and $\norm g = 1$.
\end{thm}

\begin{proof}[\emph{Hint for proof}] Let $N := \spn{(M \cup \{y\})}$ and define
$f\colon N \sto \K \colon x + \alpha y \mapsto \alpha d$ ($x \in M$, $\alpha \in \K$). \ns
\end{proof}

The most useful special case of the preceding theorem is when $M = \{\vc 0\}$ and $y \ne \vc 0$.
This tells us that the only way for $f(x)$ to vanish for every $f \in V^*$ is for $x$ to
be the zero vector.

\begin{cor}\label{HBCorI} Let $V$ be a normed linear space and $x \in V$.  If $f(x) = 0$ for
every $f \in V^*$, then $x = \vc 0$.
\end{cor}

\begin{defn} A family $\fml G(S)$ of scalar valued functions on a set~$S$
 \index{separation!of pairs of points by a family of functions}%
\df{separates} points of $S$ if for every pair $x$ and $y$ of distinct points in $S$ there exists
a function $f \in \fml G(S)$ such that $f(x) \ne f(y)$.
\end{defn}

\begin{cor}\label{HBCorII} If $V$ is a normed linear space, then $V^*$ separates points of~$V$.
\end{cor}

\begin{cor}\label{HBCorIII} Let $V$ be a normed linear space and $x \in V$.  Then there exists
$f \in V^*$ such that $\norm f = 1$ and $f(x) = \norm x$.
\end{cor}

\begin{prop} Let $V$ be a normed linear space, $\{x_1,\dots,x_n\}$ be a linearly independent
subset of $V$, and $\alpha_1, \dots, \alpha_n$ be scalars. Then there exists an $f$ in
$V^*$ such that $f(x_k) = \alpha_k$ for $1\le k \le n$.
\end{prop}

\begin{proof}[\emph{Hint for proof}] For $x = \sum_{k=1}^n \beta_k x_k$ let $g(x) =
\sum_{k=1}^n \alpha_k \beta_k$.  \ns
\end{proof}

\begin{prop}\label{C067131} Let $V$ be a normed linear space and $\vc 0 \ne x \in V$.  Then
   \[ \norm x = \max\{\abs{f(x)}\colon f \in V^* \text{ and } \norm f \le 1\}\,. \]
\end{prop}

The use of $\max$ instead of $\sup$ in the preceding proposition is intentional.  The proposition claims that there
exists an $f$ in the closed unit ball of $V^*$ such that $f(x) = \norm x$.

\begin{defn} A subset $A$ of a topological space $X$ is
 \index{dense}%
\df{dense} in $X$ if $\clo A = X$.  The space $X$ is
 \index{separable}%
\df{separable} if it contains a countable dense subset.
\end{defn}

\begin{exam} The space $\K^n$ is separable under its usual topology; that is
 \index{usual!topology!for $\K^n$}%
 \index{topology!on $\K^n$}%
 \index{K@$\K^n$!usual topology on}%
the topology it inherits from its usual norm. (See \ref{C023125}, \ref{0001503}, and \ref{xmpl_top_metric}.)
\end{exam}

\begin{exam} Given the discrete topology the real line $\R$ is not separable.
\end{exam}

\begin{prop} Let $V$ be a normed linear space.  If $V^*$ is separable, then so is~$V$.
\end{prop}






















\endinput
